\subsection{利用者と開発者の人間要因}

ソフトウェア利用者と開発者は、異なる目的、知識、行動様式を持つ。人間とコンピュータの相互作用（HCI）は、利用者が計算システムと自然にやり取りできるようにする技術と設計を扱う。利用者満足度は、利用者体験（UX）として評価されることが多い。

\subsubsection{利用者の人間要因}

利用者は、堅ろうで、安全で、直感的な画面を持ち、少ない手順で目的を達成でき、応答が速く一貫しているソフトウェアを期待する。メッセージは、成功時も失敗時も明確で完全である必要がある。可能な場合、処理中断や取り消しができることも重要である。

利用者インタフェース開発では、利用者のプロフィール、利用する入出力装置、入力方法、出力方法、障害時の許容度、性能期待を確認する。多くの場合、試作から始め、利用者の反応を取り入れて反復的に改善する。

\subsubsection{開発者の人間要因}

ソフトウェアは、開発期間より長く使われることが多い。保守する技術者が、最初に書いた技術者と同じとは限らない。したがって、コードは他者が読むことを前提に書く必要がある。意味のある文書化、適切なコーディング標準、読みやすい命名、構造化、整ったコメントは、保守性を高める。

良いコードは、書く人の都合だけでなく、後から読む人の理解を重視する。読みやすさ、単純さ、一貫したスタイル、モジュール性は、長期保守において性能と同じくらい重要な設計品質である。

\par\smallskip
\begin{center}
\centering
\IfFileExists{assets/generated_figures/ch16/fig-ch16-16.png}{\includegraphics[width=0.96\linewidth]{assets/generated_figures/ch16/fig-ch16-16.png}}{\fbox{\parbox[c][48mm][c]{0.9\linewidth}{\centering 画像生成待ち\\利用者要因と開発者要因を対比する図}}}
\par\smallskip
\refstepcounter{figure}\label{fig:ch16-16}図 \thefigure: 利用者要因と開発者要因を対比する図
\end{center}
図 \thefigure は、利用者要因と開発者要因を対比する図を視覚的に整理したものである。
\par\smallskip
